\section{Tipos de Fine-Tuning y escenarios de aplicación}

\subsection{Tres tipos de Fine-Tuning}

El ponente divide el Fine-Tuning en tres categorías, cada una con distintas necesidades de escala de datos y objetivos:

\begin{table}[H]
\centering
\caption{Comparación de tipos de Fine-Tuning}
\begin{tabular}{lccc}
\toprule
\textbf{Tipo} & \textbf{Objetivo} & \textbf{Escala de datos} & \textbf{Ejemplo} \\
\midrule
Propósito general (General Purpose) & Crear un chatbot polivalente & $>$1,000,000 muestras & ChatGPT, Claude \\
Específico de dominio (Domain Specific) & Incorporar conocimiento de dominio & miles$\sim$cientos de miles & medicina/derecho/finanzas \\
Específico de tarea (Task Specific) & Aprender una única función & cientos$\sim$miles & resumen, corrección ortográfica \\
\bottomrule
\end{tabular}
\end{table}

\begin{warningbox}{La necesidad de datos depende de múltiples factores}
Las cantidades de muestras anteriores son solo estimaciones aproximadas. La necesidad real depende de: (1) \textbf{el tamaño del modelo}: cuanto más grande es el modelo, menos muestras se necesitan; (2) \textbf{la complejidad de la tarea}: si el conocimiento del base model sobre el idioma o dominio objetivo es cercano a cero (como una lengua poco común que nunca ha visto), el efecto del fine-tuning será muy limitado, porque no hay suficientes tokens para «aprender desde cero».
\end{warningbox}

\subsection{Cuándo conviene hacer Fine-Tuning}

El ponente recomienda seguir un proceso de decisión claro:

\begin{enumerate}
    \item \textbf{Probar primero In-Context Learning / RAG}: es la opción con la barrera más baja y permite una validación rápida.
    \item Si el resultado no es satisfactorio, considerar las siguientes cuatro motivaciones para el fine-tuning:
    \begin{itemize}
        \item \textbf{Cambiar el estilo o el formato de la salida}: por ejemplo, ajustar el tono con el que el modelo genera correos electrónicos
        \item \textbf{Agregar conocimiento de dominio} (superficial): por ejemplo, enseñarle al modelo hechos sobre tu empresa
        \item \textbf{Destilación de conocimiento}: destilar capacidades del nivel de GPT-4 en un modelo más pequeño, para reducir el costo y la latencia
        \item \textbf{Mejorar la calidad en una tarea específica}: por ejemplo, en la generación de diagram, es posible superar a un frontier model en una tarea acotada
    \end{itemize}
\end{enumerate}

\begin{knowledgebox}{Motivaciones comerciales de las empresas para elegir Fine-Tuning}
Además del aspecto técnico, hay dos razones importantes no técnicas por las que las empresas eligen la ruta de modelo de código abierto más fine-tuning: (1) \textbf{control de los datos}: no quieren enviar datos sensibles a la nube mediante una API; (2) \textbf{personalización}: quieren personalizar a fondo el comportamiento, el tono y el formato de salida del modelo, para crear un modelo propio.
\end{knowledgebox}

\subsection{Resumen de la sección}

El Fine-Tuning se divide en tres niveles: propósito general, específico de dominio y específico de tarea, con una necesidad de datos que disminuye sucesivamente. En la práctica, primero se debe probar in-context learning y, solo después de confirmar que es insuficiente, considerar el fine-tuning. El valor del Fine-Tuning no está solo en la mejora técnica, sino también en dar a las empresas el control sobre el modelo y los datos.

\newpage
